\begin{statement}{A.32}{Proved}{The other symmetric genus-two quotients at every prime}
For \(a=6\) or \(-6\), put
\[
 D_a:\quad z^2=(t^2-4)(t^3-3t+a),\qquad
 Y_a=D_a\setminus\{Q_\infty\}.
\]
Let \(P_{a,\ell}\) be the geometric pro-\(\ell\) group of
\(Y_a\), based at the selected lift of the original cusp
packet of Part A. For every prime \(\ell\), every original
packet, and every permitted open \(\Delta_0\subset\Delta\),
\[
             P_{a,\ell}^{\,s_b(\Delta_0)}=I_{a,\ell}.
                                                        \tag{A.32.1}
\]
Here \(I_{a,\ell}\) is the full inertia of that selected cusp.
The assertion applies to every reflection quotient of the two
nonstable symmetric covers. Together with A.27 and A.31 it
determines all the symmetric genus-two quotient fixed groups
at every prime.
\end{statement}
\begin{proof}
\paragraph{The arithmetic and the pointed reduction.}
The actual map to the original elliptic curve is
\[
 f_a:D_a\longrightarrow E,\qquad
 x_E=\frac{t^3-3t+a/3}{4},\qquad
 y_E=\frac{z(t^2-1)}8.                                  \tag{A.32.2}
\]
Substitution, using \(a^2=36\), gives \(y_E^2=x_E^3-x_E\).
The pole order of \(x_E\) is six, so the map has degree
three and \(f_a^*O=3Q_\infty\). It is precisely the
descent of the cover in A.4, and is etale on the displayed
open curves. The complementary elliptic curve is
\[
                  F_a:\quad u^2=t^3-3t+a,
                  \qquad j(F_a)=-216.
\]
A.4 gives \(J_{D_a}\sim E\times F_a\) over a finite
number field. The non-CM assertion for \(F_a\) was proved
in A.14: for \(a=6\) its ordinary Frobenius discriminants
at five and eleven are \(-19\) and \(-35\); the other
sign is a geometric twist. An imaginary quadratic
endomorphism field could not be both of these fields.

Choose a finite number field \(L\) defining these data
and with \(\Gamma=\overline{[G_L,G_L]}\subset\Delta_0\).
To see this last inclusion, choose an open neighborhood
\(V\) of the identity in \(G_K\) with
\(V\cap\Delta\subset\Delta_0\), then take \(G_L\subset V\).
Enlarge \(L\) further to realize semistable reduction at a
place above two. Both elliptic factors have integral
\(j\), so both acquire good reduction there. Thus the
Jacobian acquires good reduction and the stable dual
graph of \(D_a\) is a tree.

We spell out why the two-tail model of A.25 still applies.
Its proof did not require the two elliptic factors to be
isogenous. The good special fibre of the original \(E\)
is supersingular in characteristic two. If \(D_a\) had
smooth stable reduction, (A.32.2) would specialize to a
separable cubic map with its only ramification a tame
point of index three. A cyclic Galois closure is
impossible by the degree-zero divisor condition modulo
three. An \(S_3\) closure would give a connected etale
double cover of the supersingular elliptic target,
also impossible. The compact-type genus-two fibre is
therefore two elliptic curves meeting in one node.

The map extends to the stable model by resolving its
rational map and contracting the rational exceptional
components, exactly as in A.25. Its two component degrees
are positive: a degree-three map on one elliptic component
would be etale and have three distinct points over the
origin, whereas the horizontal divisor \(3Q_\infty\)
and the node supply at most two. Their sum is three,
so the degrees are one and two. The divisor
\(f_a^*O=3Q_\infty\) has no vertical component, since
both component maps are nonconstant. If the cusp met
the smooth part of one component, the other component
would have no preimage of the origin, contradicting
surjectivity. Thus the cusp specializes to the node.
The pointed model inserts a rational bridge containing
the cusp, with two strictly positive arm lengths.

\paragraph{The separating derivation.}
Write \(V=V_1\oplus V_2\) for component homology and
choose symplectic bases \(a_1,b_1\) and \(a_2,b_2\).
Put \(\omega_i=[a_i,b_i]\) and
\(\Omega=\omega_1+\omega_2\).
For odd \(\ell\), apply the affine edge-twist theorem
to an ordinary point on the rational bridge. A local
Frobenius--tame commutator belongs to \(\Gamma\), acts
trivially on homology, and has first possible outer
Johnson degree two, with representative
\[
                        m_1T_1+m_2T_2,
                         \qquad m_1m_2\ne0.             \tag{A.32.3}
\]
Here the \(T_i\) are those of A.27, and \(m_i\) is
the corresponding positive arm length times a common
nonzero tame scalar. The actual packet action differs
only by an inner action. Its possible degree-one term
is \(\operatorname{ad}(w_1)\). Since it fixes the actual
cusp, \([w_1,\Omega]=0\), and the free-Lie centralizer
of \(\Omega\) gives \(w_1=0\). In degree two the same
argument makes the inner term a scalar multiple of
\(\operatorname{ad}(\Omega)\). Hence the actual
derivation is \(D_\lambda\) of A.27. This reasoning
uses only the local commutator's trivial homology
action; the whole \(\Gamma\)-action on homology need
not be trivial.

At \(\ell=2\), A.31's tangential semistable comparison
applies to this same pointed model and to every
induced packet. In its crystalline graded realization,
the identical boundary-centralizer argument gives
\[
                       N_1=0,\qquad N_2=D_\lambda.
                                                        \tag{A.32.4}
\]
The two coefficients in (A.32.3) are nonzero node
thicknesses, with their common monodromy normalization.
No comparison of a pro-two group with its prime-to-two
covering quotient is involved.

\paragraph{The non-CM factor supplies the missing invariants.}
Let \(V_E=f_a^*H_1(E)\) and let \(V_F\) be its
polarization-orthogonal complement. They are actual
\(G_L\)-subrepresentations, with
\(V=V_E\oplus V_F\), and \(V_F\simeq H_1(F_a)\).
The action of \(\Gamma\) on \(V_E\) is trivial,
because \(E\) has CM defined over \(L\).
Serre's open-image theorem for the non-CM curve
\(F_a\) gives an open image in
\(\mathrm{GL}_2(\mathbb Z_\ell)\), for every \(\ell\)
\cite[Introduction (1)--(2), Theorem 3]{A_Serre1972}.
Its commutator image is Zariski dense in
\(\mathrm{SL}(V_F)=\mathrm{Sp}(V_F)\): the Lie
algebra of that image contains
\([\mathfrak{gl}_2,\mathfrak{gl}_2]=\mathfrak{sl}_2\),
and its determinant is one. Consequently every
\(\Gamma\)-fixed vector of \(\mathbb L_d(V)\) is
fixed by the algebraic group \(\mathrm{Sp}(V_F)\),
acting as the identity on \(V_E\).

We compute the relevant intersection over any
characteristic-zero coefficient field. The component
degrees one and two permit normalizations
\[
 f_*a_1=x,\quad f_*b_1=y,\quad
 f_*a_2=x,\quad f_*b_2=2y.
\]
Adjointness gives
\[
 \begin{split}
 V_E&=\langle A,B\rangle,
       \qquad A=a_1+2a_2,\quad B=b_1+b_2,\\
 V_F&=\langle C,D\rangle,
       \qquad C=a_1-a_2,\quad D=2b_1-b_2.
 \end{split}                                             \tag{A.32.5}
\]
The two pairings are \(\langle A,B\rangle
=\langle C,D\rangle=3\); all cross pairings vanish.
Choose \(g\in\mathrm{Sp}(V_F)\) given by
\(C\mapsto2C\), \(D\mapsto D/2\), and the identity
on \(V_E\). In the \(a\)-coordinates its matrix is
\[
 M=\frac13\begin{pmatrix}5&-1\\-2&4\end{pmatrix},
\]
and its \(b\)-matrix is \((M^{-1})^T\).
Every entry of \(M\) is nonzero. Thus, for the
quadratic plane \(U=\langle\omega_1,\omega_2\rangle\),
\[
                              U\cap gU=k\Omega.         \tag{A.32.6}
\]
Indeed this is the intersection of diagonal matrices
with \(M\operatorname{Diag}M^{-1}\). If
\(DM=MD'\) with \(D,D'\) diagonal, its four entries
force all diagonal entries to be the same.
The tensor-intersection proof of A.27 then gives
\[
              \mathbb L(U)\cap\mathbb L(gU)=k\Omega.
                                                        \tag{A.32.7}
\]
The explicit matrix \(g\) is an algebraic symmetry
of the invariant space; we do not claim it is itself
a specified Galois element.

\paragraph{The prime-two invariant comparison.}
For odd \(\ell\), a leading class of a fixed element
is in \(\ker D_\lambda\) and is fixed by
\(\mathrm{Sp}(V_F)\). A.27 immediately gives the
line \(k\Omega\) when \(\lambda\ne0\); when
\(\lambda=0\), (A.32.7) gives the same conclusion.

For \(\ell=2\), let \(W_d\) be the global orbit
span of the leading class. The corresponding orbit
span in a sufficiently large finite Lie quotient
is fixed by \(\Gamma\) and has abelian global
Galois image. A.31 proves it crystalline, so
\[
                    D_{\rm cris}(W_d)\subset\ker N_2.
                                                        \tag{A.32.8}
\]
We must also transport the algebraic invariants;
non-CM global operators need not commute with the
local Galois action. Define the invariant subspace
instead as the kernel of the canonical infinitesimal
action morphism
\[
 \partial:\mathbb L_d(V)\longrightarrow
       \mathbb L_d(V)\otimes\mathfrak{sl}(V_F)^\vee,
              \qquad \partial(w)(T)=T(w).                \tag{A.32.9}
\]
This is a morphism of crystalline representations:
both summands of \(V\) are crystalline, and the
map is formed from their projections, tensors and
brackets. Exact tensoriality of \(D_{\rm cris}\)
therefore identifies its kernel with the analogous
kernel on crystalline realizations. Since
\(\mathrm{Sp}(V_F)\) is connected in characteristic
zero, its group invariants equal these Lie-algebra
invariants. Thus \(D_{\rm cris}(W_d)\) is fixed
by every element of
\(\mathrm{Sp}(D_{\rm cris}V_F)\).

Component maps and their polarization adjoints
give (A.32.5) also in crystalline homology; the
degree-two component map is invertible there
over the coefficient field. Apply the abstract
symplectic matrix in (A.32.6) within this
realization. Equations (A.32.8), (A.32.7) and
A.27 put \(D_{\rm cris}(W_d)\) on the line
\(\Omega\) in degree two. That line is the
crystalline image of the etale cusp line.
Faithfulness of the crystalline functor applied
to the quotient by that line gives
\(W_d\subset\mathbb Q_2\Omega\).
This argument never applies Fontaine monodromy
to an individual etale vector.

\paragraph{The full pro-prime conclusion.}
For every prime we have proved that a nonidentity
fixed element has first nonzero degree two and
leading class on the actual cusp line. In the
integral free Lie algebra that cusp class is
primitive: it is the sum of the two standard
handle commutators. Its coefficient is therefore
in \(\mathbb Z_\ell\), including at two and three.
Subtract the corresponding actual cusp power.
The remainder is fixed and has strictly larger
first degree, which the proved restriction
excludes. Separation of the lower-central
filtration gives the identity. The reverse
inclusion follows from the trivial cyclotomic
character on \(\Delta_0\). This proves (A.32.1).
\end{proof}
